\chapter{Operazioni sulle funzioni e Paradigma Funzionale}
\section{Paradigma Funzionale}
Fino a ora abbiamo utilizzato un approccio \textit{diretto} per la valutazione della calcolabilità di una funzione. Questo approccio presenta
svariati limiti, primo dei quali è la necessità di dimostrare come calcolabile una singola funzione alla volta, senza poter
trarre conclusioni  di carattere generale o di caratterizzare insiemi o classi di funzioni.

Per avere un approccio più generale e sistematico alla valutazione della calcolabilità delle funzioni è necessario dunque
approcciare il problema in maniera diversa, utilizzando un paradigma diverso che sfrutti\footnote{Il passaggio dall'approccio diretto a uno di carattere più generale per la valutazione della calcolabilità di una funzione è assimilabile al passaggio dal paradigma imperativo a quello funzionale.}
le proprietà e le operazioni delle funzioni stesse per valutarne la calcolabilità.

\section{Operazioni fondamentali e Classi PRC}

Con \textbf{operazioni sulle funzioni} facciamo riferimento a operazioni algebriche definite nell'insieme delle funzioni, ovvero operazioni che hanno sia come operando che come risultato delle funzioni.
Definiamo dunque le due operazioni fondamentali sulle funzioni: la \textbf{composizione} e la \textbf{ricorsione primitiva}.

\subsection{Composizione}

\dfn{Composizione di Funzioni}{\label{Composizione di funzioni Def}
  Sia \(f\) una funzione \(k\text{-aria}\) e \( g_1, g_2, \cdots,g_k\) funzioni \(n\text{-arie}\). Allora la funzione \(h\) \(n\text{-aria}\) è \textbf{composizione} di \(f, g_1, g_2, \cdots,g_k\) se:
  \begin{equation}
    h(x_1,\cdots,x_n)=f(g_1(x_1,\cdots,x_n),\cdots,g_k(x_1,\cdots,x_n))
  \end{equation}
}

È importante notare che la composizione di funzioni è ben definita anche per \textit{funzioni parziali}. \\
Infatti si avrà che, data \((x_1,\cdots,x_n) \in \mbN\):

\begin{equation}
  h(x_1,\cdots,x_n)\uparrow \iff \exists i \in \cbrakets{1,\cdots,k}: g_i(x_1,\cdots,x_n)\uparrow \vee f(g_1(x_1,\cdots,x_n),\cdots,g_k(x_1,\cdots,x_n))\uparrow
\end{equation}

Con ragionamento simile è anche possibile dimostrare che la composizione di \textbf{funzioni parzialmente calcolabili} è necessariamente \textbf{parzialmente calcolabile}.

\newpage

\begin{halfframedbox}{red!75!black}{Teorema di Composizione}\label{Composizione Teorema}
  \textbf{Enunciato.} Sia \(h\) ricavata per composizione di \(f, g_1, g_2, \cdots,g_k\) \textbf{parzialmente calcolabili}, allora \(h\) è \textbf{parzialmente calcolabile}.

  \begin{center}
    \textcolor{red!75!black}{\hdashrule[0.5ex]{18cm}{1pt}{3mm 3pt 1mm 2pt}}
  \end{center}

  \textbf{Dimostrazione.} Per definizione di \textbf{composizione} si ha \(h=\Psi_{\mcP}^{(n)}\), dove \(\mcP\) è:
  \begin{lstlisting}[language=ini, label={Teorema compozione},  caption=Dimostrazione Teorema di Composizione]
    Z1 <- (*\(g_1(x_1,\cdots,x_n)\)*)
    (*\(\vdots\)*)
    Zk <- (*\(g_k(x_1,\cdots,x_n)\)*)
    Y <- (*\(f(Z_1,\cdots,Z_k)\)*)
  \end{lstlisting}
\end{halfframedbox}

Inoltre come notato in precedenza il programma non arriva a terminazione \textit{se e solo se} uno dei programmi \(g_i\) o \(f\) non termina, dunque:
\begin{equation}
   h \text{ non totale } \implies \exists i \in \cbrakets{1,\cdots,k}: g_i \text{ non totale} \lor f \text{ non totale}
\end{equation}
Non è necessariamente vero il viceversa però. Infatti se il codominio di una funzione totale \(f\) che compone \(h\) è un sottoinsieme del dominio di un'altra parziale \(g\) che compone a sinistra\footnote{In notazione algebrica \(f \circ g\) equivale a \(f(g(x))\), dunque è l'operando sinistro a poter essere parziale, non essendo la composizione commutativa.},
la funzione composta sarà totale, poiché quella totale escluderà i casi gli input di non terminazione dalla prima.
Vediamo qualche esempio al riguardo per chiarire il concetto.

\begin{generalbox}[colframe=azure-gradient-3!90!black]{Esempi}
  \begin{itemize}
    \item Sia \(h_1(x)=2x\). Possiamo considerare \(h\) come funzione composta da \(f_1(x_1, x_2)=x_1 \cdot x_2\), \(g_1(x)=2\) e \(g_2(x)=x\), o alternativamente, notando che \(h_1(x)=2x=x+x\) ovvero composizione di \(f_2(x_1,x_2)=x_1+x_2\) e \(g_2(x)=x\), dunque \(h_1\) è calcolabile.
    \item Sia \(h_2(x)=4x^2\). Possiamo considerare \(h_2\) come funzione composta dalle funzioni \(f_1\) e \(h_1\) sopracitate, infatti \(h_2(x)=4x^2=2x\cdot 2x\), dunque \(h_2\) è calcolabile.
    \item Sia \(h_3(x)=4x^2-2x\). Si ha \(h_3(x)\downarrow \forall x \in \mbN\), dato che \(\forall x \in \mbN, 4x^2 \geq 2x\). Abbiamo dunque che la funzione \(h_3\) è totale. Ma \(h_3(x)=f_3(h_2(x),h_1(x))\), dove \(f_3(x_1,x_2)=x_1-x_2\) \textbf{non è totale}\footnote{In caso di dubbi controllare qui~\ref{Funzione Sottrazione}}. Abbiamo dunque che l'implicazione non vale in entrambi i sensi come precedentemente detto.
  \end{itemize}
\end{generalbox}

\subsection{Ricorsione Primitiva}

L'operazione di \textbf{ricorsione primitiva} è un'operazione che permette di definire una funzione in maniera ricorsiva, ovvero definendo un caso base e un passo ricorsivo, formalizzando matematicamente il concetto di algoritmo ricorsivo.
Nonostante sia definibile anche per le \textbf{funzioni parziali}, noi tratteremo solo il caso delle \textbf{funzioni totali}.

Per una funzione \(h\) unaria diremo che è ricavata per \textbf{ricorsione primitiva} da \(k\in\mbN\) e \(g\) funzione binaria \textbf{totale} se:
\begin{equation}
  \begin{cases}
    h(0)=k\\
    h(t+1)=g(t,h(t)), & \forall t\in\mbN^*
  \end{cases}
\end{equation}

Nonostante affronteremo prevalentemente funzioni unarie, è importante dare una definizione generale per qualsiasi funzione.

\dfn{Ricorsione Primitiva}{\label{Ricorsione Primitiva Def}
Una funzione \(h\) \(n-aria\) si ricava per ricorsione primitiva da \(f\) \((n-1)\text{-aria}\) e da \(g\) \((n+1)\text{-aria}\), se \(\forall x_1,\cdots,x_{n-1} \in \mbN\) vale:
\begin{equation}
  \begin{cases}
    h(x_1,\cdots,x_{n-1}, 0)=f(x_1,\cdots,x_{n-1})\\
    h(x_1,\cdots,x_{n-1}, t+1)=g(x_1,\cdots,x_{n-1},t,h(x_1,\cdots,x_{n-1}, t)), & \forall t\in \mbN^*
  \end{cases}
\end{equation}
}

\newpage

\begin{halfframedbox}{red!75!black}{Teorema di Ricorsione primitiva}\label{Ricorsione Primitiva Teorema}
  \textbf{Enunciato.} Sia \(h\) ricavata per ricorsione primitiva da \(f\) e \(g\) calcolabili, allora \(h\) è calcolabile.

  \begin{center}
    \textcolor{red!75!black}{\hdashrule[0.5ex]{18cm}{1pt}{3mm 3pt 1mm 2pt}}
  \end{center}

  \textbf{Dimostrazione.} Per definizione di \textbf{ricorsione primitiva} si ha \(h=\Psi_{\mcP}^{(n)}\), dove \(\mcP\) è:
  \begin{lstlisting}[language=ini, autogobble=false, label={Teorema ricorsione primitiva},  caption=Dimostrazione Teorema di Ricorsione Primitiva]
      Y <- (*\(f(X_1,\cdots,X_{n-1})\)*)
  [A] IF Xn == 0 GOTO E
      Y <- (*\(g(X_1,\cdots,X_{n-1},Y,Z)\)*)
      Z <- Z + 1
      Xn <- Xn - 1
      GOTO A
  \end{lstlisting}
\end{halfframedbox}

Questo programma calcola esattamente la funzione \(h\) poiché, se ci trovassimo nel primo caso della definizione~\ref{Ricorsione Primitiva Def} l'output sarebbe dato dalle linee \lstinline[language=ini]{1--2}, altrimenti, se ci trovassimo nel secondo caso, le linee \lstinline[language=ini]{3--6} calcolerebbero ricorsivamente l'output di \(h\), dimostrata così calcolabile.

\subsection{Classi PRC e Funzioni Ricorsive Primitive}

Le due operazioni appena definite sono alla base del paradigma funzionale e definiscono, insieme ad alcune specifiche funzioni, una classe di funzioni chiamate \textbf{PRC}.

\dfn{Classi PRC di funzioni}{\label{Classi PRC di funzioni Def}
  Un insieme di funzioni totali \(\mcC\) è detto \textbf{PRC} se è chiuso rispetto alle operazioni di composizione e ricorsione primitiva e contiene:
  \begin{itemize}
    \item La funzione nulla \(n(x)=0\);
    \item La funzione successore \(s(x)=x+1\);
    \item Le funzioni proiezione \(u_i^{(n)}(x_1,\cdots,x_n)=x_i, \forall (n\geq 1 \land 1\leq i \leq n )\).
  \end{itemize}
}

Dato che le funzioni iniziali\footnote{Si noti che le funzioni proiezioni sono un \textbf{insieme infinito} e non \textbf{una sola funzione}. Infatti ne esiste una \(\forall (n,i)\) con \( n \in \mbN\) e \( i \in \cbrakets{1,\cdots,n}\).} sono facilmente calcolabili e visti i risultati dei teoremi~\ref{Teorema compozione} e~\ref{Teorema ricorsione primitiva}, possiamo notare che l'insieme delle \textbf{funzioni calcolabili} è una \textbf{classe PRC}.

\dfn{Funzioni Ricorsive Primitive}{\label{Funzioni Ricorsive Primitive Def}
  Una funzione è detta \textbf{ricorsiva primitiva} se è ottenibile dalle funzioni iniziali mediante un numero finito di applicazioni di composizione e ricorsione primitiva.
}

È importante notare che, dalla definizione di \textbf{classe PRC} e dalle proprietà degli \textbf{insiemi chiusi}, l'insieme delle \textbf{funzioni ricorsive primitive} è contenuto in tutte le \textbf{classi PRC}, ed è dunque in particolare la più piccola \textbf{classe PRC}.

In particolare, essendo \textbf{l'insieme delle funzioni calcolabili} una \textbf{classe PRC}, si ha che \textbf{ogni funzione ricorsiva primitiva è calcolabile}.

\newpage
Andiamo ora a elencare una serie di esempi significativi di funzioni ricorsive primitive, che verranno poi usate sia per definire nuove funzioni che per dimostrare l'appartenenza di altre funzioni alla stessa classe.

\begin{generalbox}[colframe=azure-gradient-3!90!black]{Esempi}
  \begin{enumerate}
    \item Consideriamo la funzione \(a(x_1,x_2) = x_1+x_2\). Possiamo vedere che questa è \textbf{ricorsiva primitiva} da:
    \begin{itemize}
      \item []
      \begin{equation}
        \begin{cases}
          a(x_1,0)=x_1=u_1^{(1)}(x)\\
          a(x_1,t+1)=x+t+1=a(x_1,t)+1=s(a(x_1,t))=g(t, a(x,t),x), & \forall t \in \mbN^*
        \end{cases}
      \end{equation}
      Dove \(g(x,y,z)=s(y)=s(u_2^{(3)}(x,y,z))\) è ottenuta per composizione di funzioni iniziali, ed è dunque \textbf{ricorsiva primitiva}.
    \end{itemize}
    \item Consideriamo la funzione \(m(x,y)=x\cdot y\). Possiamo vedere che questa è \textbf{ricorsiva primitiva} da:
    \begin{itemize}
     \item []
      \begin{equation}
        \begin{cases}
          m(x,0)=0=n(x)\\
          m(x, t+1)=m(x, t)+x=a(m(x,t),x)=g'(t, m(x, t),x), & \forall t \in \mbN^*
        \end{cases}
      \end{equation}
      Dove \(g'(x,y,z)=a(u_2^{(3)}(x,y,z),u_3^{(3)}(x,y,z))\) è ottenuta per composizione di funzioni iniziali, ed è dunque \textbf{ricorsiva primitiva}.
    \end{itemize}
    \item Consideriamo la funzione \textbf{fattoriale \(x!\)}, definita come: % chkTeX 40
    \begin{itemize}
      \item []
      \begin{equation}
        x! =\begin{cases}
              1, & x=0\\
              x\cdot (x-1)!,&  x>0
            \end{cases}
      \end{equation}
      Possiamo vedere che questa è \textbf{ricorsiva primitiva} da:
      \begin{equation}
        \begin{cases}
          0! =1=s(0)\\
          (t+1)! =(t+1)\cdot t! =m(t+1, t!)=g(t+1,t!), & \forall t \in \mbN^*
        \end{cases}
      \end{equation}
    \end{itemize}
    \item Consideriamo la funzione \textbf{precedente}:
    \begin{itemize}
     \item []
     \begin{equation}
        p(x)=\begin{cases}
              0, &  x=0\\
              x-1, & x>0
            \end{cases}
      \end{equation}
      Possiamo vedere che questa è \textbf{ricorsiva primitiva} da:
      \begin{equation}
        \begin{cases}
          p(0)=0=n(x)\\
          p(t+1)=t=u_1^{(2)}(t,p(t))=g(t, p(t)), & \forall t \in \mbN^*
        \end{cases}
      \end{equation}
    \end{itemize}
    \item Consideriamo ora un'alternativa alla funzione Sottrazione:
    \begin{itemize}
     \item []
      \begin{equation}
        x\dotminus y =  \begin{cases}
                          0, & x<y \\
                          x-y, & x\geq y
                        \end{cases}
      \end{equation}
      Possiamo vedere che questa è \textbf{ricorsiva primitiva} da:
      \begin{equation}
        \begin{cases}
          x\dotminus 0=x=u_1^{(1)}(x)\\
          x\dotminus (t+1)=p(x\dotminus t)=g(x, x\dotminus t,t), & \forall t \in \mbN^*
        \end{cases}
      \end{equation}
      Dove \(g(x,y,z)=p(u_2^{(3)}(x,y,z))\) è ottenuta per composizione di funzioni iniziali, ed è \textbf{ricorsiva primitiva}\footnote{Per dimostrare che una funzione è \textbf{ricorsiva primitiva} è sufficiente dimostrare che è ottenuta da \textbf{funzioni ricorsive primitive}. Da ora in avanti non ci si ricondurrà più a funzioni iniziali per le dimostrazioni.}.
    \end{itemize}
    \item Consideriamo ora il predicato\footnote{Questo predicato è detto \textit{rilevatore di zeri}, è applicato a un predicato può svolgere il ruolo della \textit{negazione logica}} \(\alpha(x)\iff x=0\), esprimibile anche come:
    \begin{itemize}
     \item []
     \begin{equation}
        \alpha(x)=\begin{cases}
                    1, & x=0\\
                    0, & x>0
                  \end{cases} = 1\dotminus x
      \end{equation}
      Questo predicato è \textbf{ricorsivo primitivo}, essendo composizione di funzioni dimostrate come tali.
    \end{itemize}
    \item Consideriamo ora il predicato \(\beta(x,y)\iff x=y\), esprimibile anche come\footnote{Funzione \textit{Sottrazione in Valore Assoluto}:~\ref{Funzione Sottrazione Valore Assoluto Eq}}:
    \begin{itemize}
     \item []
     \begin{equation}
        \beta(x,y)=\begin{cases}
                    1, & x=y\\
                    0, & x\neq y
                  \end{cases} = \alpha(\abs{x-y})
      \end{equation}
      Questo predicato è ricorsivo primitivo, essendo composizione di funzioni dimostrate come tali\footnote{La funzione \(f(x,y)=\abs{x-y}\) è definita e dimostrata come ricorsiva primitiva nell'esercizio~\ref{Funzione Sottrazione Valore Assoluto Eq}. Esercizi consigliati:~\ref{Funzione Sottrazione Valore Assoluto Eq},~\ref{Predicato Minore Uguale Eq}}.
    \end{itemize}
  \end{enumerate}
\end{generalbox}

\newpage
\section{Proprietà di chiusura delle Classi PRC}

Definite le classi PRC e, notato che l'insieme delle funzioni calcolabili ne è un esempio, è naturale chiedersi se esse siano chiuse rispetto ad altre operazioni oltre a quelle che le definiscono.

Proprietà di chiusura di questo tipo sono importanti per poter caratterizzare le classi PRC e per poter dimostrare l'appartenenza di una funzione a una di esse, e dunque la sua calcolabilità.

Andiamo dunque a elencare e dimostrare le proprietà di chiusura delle classi PRC. % chktex 13

\subsection{Chiusura per congiunzione, disgiunzione e negazione}

\begin{halfframedbox}{red!75!black}{Teorema di Chiusura per congiunzione\(\text{,}\) disgiunzione e negazione}\label{Chiusure per congiunzione, disgiunzione e negazione Teorema}
  \textbf{Enunciato.} Siano \(p\) e \(q\) predicati n-ari appartenenti a \(\mcC\) \textbf{classe PRC}, allora anche i predicati \(p \land q\), \(p \lor q\), \(\neg p\), \(\neg q\) appartengono a \(\mcC\).

  \begin{center}
    \textcolor{red!75!black}{\hdashrule[0.5ex]{18cm}{1pt}{3mm 3pt 1mm 2pt}}
  \end{center}

  \textbf{Dimostrazione. }
  \begin{itemize}
    \item \((p \land q)(x_1,\cdots,x_n) = p(x_1,\cdots,x_n) \cdot q(x_1,\cdots,x_n)\), è dunque \textbf{ricorsivo primitivo} da composizione di funzioni ricorsive primitive.
    \item \((\neg p)(x_1,\cdots,x_n)=\alpha(p(x_1,\cdots,x_n))\), è dunque \textbf{ricorsivo primitivo} da composizione di funzioni ricorsive primitive, vale il ragionamento analogo per \(\neg q\).
    \item \((p \lor q)(x_1,\cdots,x_n) \iff \neg(\neg p(x_1,\cdots,x_n) \land \neg q(x_1,\cdots,x_n))\), è dunque \textbf{ricorsivo primitivo} dalle leggi di De Morgan e dal ragionamento precedente.
   \end{itemize}
\end{halfframedbox}

Avendo dimostrato che i predicati \(x \leq y\) e \(x=y\) sono \textbf{ricorsivi primitivi} e utilizzando il teorema appena dimostrato, sappiamo che sono ricorsivi primitivi anche i \textbf{predicati} \(x<y\iff x\leq y \land \neg (y= x)\), \(x\geq y\iff\neg(x<y)\) e \(x>y\), dalle precedenti.

\subsection{Definizione per Casi}

Andiamo ora a individuare un ulteriore metodo per definire funzioni che, come vedremo, è una proprietà di chiusura per le classi PRC, ovvero la \textbf{definizione per casi}.

\dfn{Definizione per Casi}{
  Siano \(g\) e \(h\) funzioni \(n\text{-arie}\) e \(p\) predicato \(n\text{-ario}\) allora la funzione \(n\text{-aria}\) \(f\) è \textbf{definita per composizione} se definita come:
  \begin{equation}
    f(x_1,\cdots,x_n)=\begin{cases}
                    g(x_1,\cdots,x_n), & p(x_1,\cdots,x_n)\\
                    h(x_1,\cdots,x_n), & \neg p(x_1,\cdots,x_n)
                  \end{cases}
                  \label{Definizione per Casi Def}
  \end{equation}
}

In altre parole una funzione \textbf{definita per casi} è una funzione la cui immagine dipende da un predicato, ovvero da una condizione che può essere vera o falsa riguardante le incognite della funzione stessa.
Possiamo trovare una proprietà di chiusura per le \textbf{Classi PRC} per la \textbf{definizione per casi} simile a quelle viste per composizione e ricorsione primitiva.

\begin{halfframedbox}{red!75!black}{Teorema di Definizione per Casi}\label{Definizione per Casi Teorema}
  \textbf{Enunciato.} Siano g e h funzioni \(n\text{-arie}\) in \(\mcC\) \textbf{classe PRC} e \(p\) predicato \(n\text{-ario}\) in \(\mcC\), allora a \(\mcC\) appartiene anche la funzione \(n\text{-aria}\) \(f\) \textbf{definita per casi} da \(g,h\) e \(p\).

  \begin{center}
    \textcolor{red!75!black}{\hdashrule[0.5ex]{18cm}{1pt}{3mm 3pt 1mm 2pt}}
  \end{center}

  \textbf{Dimostrazione.} Si ha:
  \begin{equation}
    f(x_1,\cdots,x_n)=g(x_1,\cdots,x_n)p(x_1,\cdots,x_n)+h(x_1,\cdots,x_n)(\neg p)(x_1,\cdots,x_n)
  \end{equation}
  Per la proprietà di chiusura per \textbf{composizione} e \textbf{ricorsione primitiva} dimostrata in precedenza, \(f\in\mcC\)
\end{halfframedbox}

È possibile generalizzare la definizione di \textbf{definizione per casi} a più casi, ovvero a più funzioni e più predicati, la cui proprietà di chiusura è  a sua volta dimostrabile\footnote{Dimostrazione alternativa:~\ref{Corollario al Teorema di Definizione per Casi AltDim}. Esercizi consigliati:~\ref{Funzione iterata n-esima di f Eq}.}.

\newpage

\begin{halfframedbox}{red!75!black}{Corollario al Teorema di Definizione per Casi}\label{Corollario al Teorema di Definizione per Casi}
  \textbf{Enunciato.} Se \(g_1,\cdots,g_k,h,p_1,\cdots,p_k\) sono funzioni e predicati n-ari in \(\mcC\) \textbf{classe PRC}, allora a \(\mcC\) appartiene anche la funzione \(n\text{-aria}\) \(f\) definita come:
  \begin{equation}
    f(x_1,\cdots,x_n)=\begin{cases}
                    g_1(x_1,\cdots,x_n), & p_1(x_1,\cdots,x_n)\\
                    \vdots \\
                    g_k(x_1,\cdots,x_n), & p_k(x_1,\cdots,x_n)\\
                    h(x_1,\cdots,x_n), & \text{altrimenti}
                    \end{cases}
  \end{equation}

  \begin{center}
    \textcolor{red!75!black}{\hdashrule[0.5ex]{18cm}{1pt}{3mm 3pt 1mm 2pt}}
  \end{center}

  \textbf{Dimostrazione.} Si può dimostrare che \(f\) è \textbf{ricorsiva primitiva} per \textbf{induzione} su \(k\). Il \textbf{caso base} è rappresentato dal teorema precedente.
  Per quando riguarda il \textbf{passo induttivo} \(k\implies k+1\), consideriamo la funzione \(f'\) definita come:
    \begin{equation}
      f'(x_1,\cdots,x_n)=\begin{cases}
                      g_k(x_1,\cdots,x_n), & p_k(x_1,\cdots,x_n)\\
                      h(x_1,\cdots,x_n), & \text{altrimenti}
                      \end{cases}
    \end{equation}
    Per il teorema precedente \(f' \in \mcC\).
    A questo punto possiamo riscrivere \(f\) come:
    \begin{equation}
      f(x_1,\cdots,x_n)=\begin{cases}
                      g_1(x_1,\cdots,x_n), & p_1(x_1,\cdots,x_n)\\
                      \vdots\\
                      g_{k-1}(x_1,\cdots,x_n), & p_{k-1}(x_1,\cdots,x_n)\\
                      f'(x_1,\cdots,x_n), & \text{altrimenti}
                      \end{cases}
    \end{equation}
    Per ipotesi di induzione, essendo ora \(f\) definita da \(k\) funzioni appartenenti a \(\mcC\) e \(k\) predicati appartenenti a \(\mcC\), \(f \in \mcC\).
  \end{halfframedbox}

\subsection{Operatori di Sommatoria e Produttoria e Quantificatori Limitati}

Come abbiamo già visto, la combinazione di funzioni \(f\in\mcC\) \textbf{classe PRC}, mediante somma e prodotto appartiene ancora \(\mcC\).
Possiamo generalizzare questo processo di composizione mediante gli operatori di \textbf{sommatoria e produttoria}.

\dfn{Sommatoria e Produttoria di funzioni}{\label{Sommatoria e Produttoria di funzioni Def}
  Sia \(f\) una funzione \((n+1)\text{-aria}\), allora definiamo come \textbf{operatori di sommatoria e produttoria} le funzioni:
  \begin{equation}
    \begin{split}
      \sigma_f(x_1,\cdots,x_n,y)=\sum_{i=0}^{y}f(x_1,\cdots,x_n,i)\\
      \pi_f(x_1,\cdots,x_n,y)=\prod_{i=0}^{y}f(x_1,\cdots,x_n,i)
    \end{split}
  \end{equation}
}

Come le operazioni precedenti, anche per gli \textbf{operatori di sommatoria e produttoria} è possibile dimostrare una proprietà di chiusura per le \textbf{classi PRC}.

\begin{halfframedbox}{red!75!black}{Teorema di chiusura per Sommatoria e Produttoria delle classi PRC}\label{Teorema di chiusura per Sommatoria e Produttoria delle classi PRC}
  \textbf{Enunciato. } Sia \(f\) una funzione \((n+1)\text{-aria}\) in \(\mcC\) \textbf{classe PRC}, allora appartengono a \(\mcC\) anche le funzioni \(\sigma_f\) e \(\pi_f\).

  \begin{center}
    \textcolor{red!75!black}{\hdashrule[0.5ex]{18cm}{1pt}{3mm 3pt 1mm 2pt}}
  \end{center}

  \textbf{Dimostrazione. } Dimostriamo per \textbf{ricorsione primitiva}. Si ha:
  \begin{equation}
    \begin{cases}
      \sigma_f(x_1,\cdots,x_n,0)=f(x_1,\cdots,x_n,0)\\
      \sigma_f(x_1,\cdots,x_n,t+1)=\sigma_f(x_1,\cdots,x_n,t)+f(x_1,\cdots,x_n,t+1), \forall t \in \mbN^*
    \end{cases}
  \end{equation}
  Si può dunque vedere che \(\sigma_f\) è \textbf{ricorsiva primitiva} da composizione di funzioni iniziali.
  Per ragionamento analogo si dimostra anche che \(\pi_f\) è \textbf{ricorsiva primitiva}.
\end{halfframedbox}

Gli \textbf{operatori di sommatoria e produttoria} sono utili per dimostrare come appartenente a \textbf{classi PRC} ulteriori costrutti matematici.
Un primo esempio sono i \textbf{quantificatori limitati esistenziale e universale}.

\newpage

\dfn{Quantificatori limitati}{\label{Quantificatori limitati Def}
  Sia \(p\) un predicato \((n+1)\text{-ario}\), allora definiamo i \textbf{quantificatori limitati esistenziale e universale} \(E_p\) e \(U_p\) come:
  \begin{equation}
    \begin{split}
    E_p(x_1,\cdots,x_n,y)\iff \exists t \leq y: p(x_1,\cdots,x_n,t)\\
    U_p(x_1,\cdots,x_n,y)\iff \forall t \leq y: p(x_1,\cdots,x_n,t)
    \end{split}
  \end{equation}
}

È fondamentale definire questi quantificatori \textbf{limitati superiormente}, poiché altrimenti non sarebbe possibile dimostrare la proprietà di chiusura delle \textbf{classi PRC} a loro relativa.

\begin{halfframedbox}{red!75!black}{Teorema di chiusura per quantificatori limitati delle classi PRC}\label{Teorema di chiusura per quantificatori limitati delle classi PRC}
  \textbf{Enunciato. } Sia \(p\) una funzione \((n+1)\text{-ario}\) in \(\mcC\) \textbf{classe PRC}, allora appartengono a \(\mcC\) anche i predicati \(E_p\) e \(U_p\) \((n+1)\text{-ari}\) definiti come:
  \begin{equation}
    E_p(x_1,\cdots,x_n,y)\iff \exists t \leq y: p(x_1,\cdots,x_n,t)
    U_p(x_1,\cdots,x_n,y)\iff \forall t \leq y: p(x_1,\cdots,x_n,t)
  \end{equation}

  \begin{center}
    \textcolor{red!75!black}{\hdashrule[0.5ex]{18cm}{1pt}{3mm 3pt 1mm 2pt}}
  \end{center}

  \textbf{Dimostrazione. }  Per \(E_p\) si ha che:
  \begin{equation}
    E_p(x_1,\cdots,x_n,y)\iff\sum_{t=0}^y p(x_1,\cdots,x_n,y)> 0
  \end{equation}
  poiché la sommatoria darà \(0\) nel caso non esista un \(t\) tale che \(p(x_1,\cdots,x_n,t)\), e viceversa un valore positivo nel caso esista.
  Per \(U_p\) si ha che:
  \begin{equation}
    U_p(x_1,\cdots,x_n,y)=\prod_{t=0}^y p(x_1,\cdots,x_n,y)
  \end{equation}
  poiché in caso esista un predicato che restituisce \(0\) la produttoria verrebbe annullata.
\end{halfframedbox}

\begin{generalbox}[colframe=azure-gradient-3!90!black]{Esempi}
  Grazie a questi quantificatori, è possibile caratterizzare altri predicati aritmetici come \textbf{ricorsive primitive}.
  \begin{enumerate}
    \item Predicato di divisibilità:
    \begin{itemize}
      \item [] \begin{equation}
        x|y \iff \exists z \leq y: x\cdot z = y
      \end{equation}
      È possibile in questo caso \textbf{limitare superiormente} la ricerca di \(z\) poiché per far si che \(x|y\) sia vero, \(z\) deve essere minore o uguale a \(y\).
    \end{itemize}
    \item Predicato numero primo:
    \begin{itemize}
      \item [] \begin{equation}
        \text{Primo}(x) \iff x>1 \land \forall z \leq x: (\neg(z|x) \lor z=1 \lor z=x)
      \end{equation}
      In questo modo esprimiamo proprio che \(x\) è primo \textit{se e solo se} ha come divisori solo \(1\) e se stesso.
      \end{itemize}
  \end{enumerate}
\end{generalbox}

\subsection{Minimalizzazione limitata}
Un ulteriore operazione per cui è dimostrabile una proprietà di chiusura delle \textbf{classi PRC} è la \textbf{minimalizzazione limitata}.
\dfn{Minimalizzazione limitata}{\label{Minimalizzazione limitata Def}
  Sia \(p\) un predicato \((n+1)\text{-ario}\), allora la funzione \((n+1)\text{-aria}\) \(\min_{t\leq y}\) è definita come:
  \begin{equation}
    \min_{t\leq y}(p(x_1,\cdots,x_n,y))=\begin{cases}
                                  \min(Z_y)=\cbrakets{t\leq y| p(x_1,\cdots,x_n,t)}, & \exists t \leq y: p(x_1,\cdots,x_n,t)\\
                                  0, & \text{altrimenti}
                                \end{cases}
  \end{equation}
}

Dove \(\min(Z_y)\) è il minimo elemento dell'insieme \(Z_y\), insieme dei valori \(t\leq\) y per cui \(p(x_1,\cdots,x_n,t)\) è vero.
Questa funzione dunque, dato un predicato \(p\) e un \textbf{limite superiore} \(y\) restituisce, se esiste\footnote{È importante specificare nella definizione di minimalizzazione limitata i due casi possibili tramite il quantificatore esistenziale, poiché il \textbf{non esiste minimo dell'insieme vuoto}.}, il minimo \(t \in \cbrakets{0, \cdots, y}\) per cui il predicato \(p\) è vero, altrimenti restituisce \(0\).

\newpage

\begin{halfframedbox} {red!75!black}{Teorema di chiusura per minimalizzazione limitata delle classi PRC}\label{Teorema di chiusura minimalizzazione limitata delle classi PRC}
  \textbf{Enunciato. } Sia \(p\) un predicato \((n+1)\text{-ario}\) in \(\mcC\) \textbf{classe PRC}, allora appartiene a \(\mcC\) anche la funzione \((n+1)\text{-aria}\) \(\min_{t\leq y}\) definita per minimalizzazione limitata da \(p\).

  \begin{center}
    \textcolor{red!75!black}{\hdashrule[0.5ex]{18cm}{1pt}{3mm 3pt 1mm 2pt}}
  \end{center}

  \textbf{Dimostrazione. } Consideriamo la seguente funzione \(g\):
  \begin{equation}
    g(x_1,\cdots,x_n,y)=\sum_{u=0}^y\prod_{t=0}^{u}\alpha(p(x_1,\cdots,x_n,t))
  \end{equation}
  Analizziamo \(g\) dall'interno. La funzione \(\alpha\) agisce sui predicati come negazione. Dunque si avrà:
  \begin{equation}
    \alpha(p(x_1,\cdots,x_n,t))= \begin{cases}
                                1, & \neg p(x_1,\cdots,x_n,t)\\
                                0, & p(x_1,\cdots,x_n,t)
                              \end{cases}
  \end{equation}
  Inevitabilmente si avrà che la produttoria in \(g\) si annullerà nel caso in cui \(p(x_1,\cdots,x_n,t)\) sia vera per qualche \(t\leq u\), ovvero:
  \begin{equation}
    \prod_{t=0}^{u}\alpha(p(x_1,\cdots,x_n,t))=1\iff\forall t \leq u, \neg p(x_1,\cdots,x_n,t)
  \end{equation}

  Di conseguenza, iterando su \(u\), otterremo che le varie \textbf{produttorie} varranno \(1\) fino a quando non si avrà il primo \(u=t_0\) tale che \(p(x_1,\cdots,x_n,u)\) risulti vera, e sarà nulla da quel momento in poi, essendoci l'iterazione \(t_0\)-esima ad annullarla.

  Dovendo infine la \textbf{sommatoria esterna} fare esattamente \(t_0\) passi prima di raggiungere tale \(u\) minimo, si avrà che \(g(x_1,\cdots,x_n,y)=t_0\), che abbiamo dimostrato essere il \textbf{minimo cercato}.

  Si avrà dunque che la funzione \(g=\text{min} Z_y\in\mcC\) essendo composizione di funzioni appartenenti a \(\mcC\). Di conseguenza anche la funzione \(\min_{t\leq y}\in\mcC\), poiché \textbf{definita per casi}.
\end{halfframedbox}

\begin{generalbox}[colframe=azure-gradient-3!90!black]{Esempi}
  Grazie alla minimalizzazione limitata è possibile definire altri predicati aritmetici come \textbf{ricorsive primitive}.
  \begin{enumerate}
    \item \textbf{Divisione con resto \(\left\lfloor \frac{x}{y} \right\rfloor\)}. Questo predicato è definibile come:
    \begin{itemize}
      \item []
      \begin{equation}
        \left\lfloor \frac{x}{y} \right\rfloor = \min_{t\leq x} ((t+1)y > x)
      \end{equation}
      Infatti il minimo \(t\) per cui \((t+1)y > x\) è proprio il quoziente cercato, poiché per \(t\) più piccoli il prodotto con \(y\) sarà minore di \(x\) e dunque non potrà essere risultato della divisione.
    \end{itemize}
    \item \textbf{Operatore di resto \(x \mod y\)}. Questo predicato è definibile come:
    \begin{itemize}
      \item []
      \begin{equation}
        x \mod y = x \dotminus \left\lfloor \frac{x}{y} \right\rfloor y
      \end{equation}
    \end{itemize}
    \item \textbf{\(N\text{-esimo}\) primo.} Definiamo \(P_0=0\) e sia per \(n>0, P_n\) l'\(n\text{-esimo}\) numero primo. Dimostriamo che questa funzione è \textbf{ricorsiva primitiva}.
    Consideriamo la funzione definita come:
    \begin{itemize}
     \item []
     \begin{equation}\label{Definizione Ricorsiva Primitiva N-esimo Primo}
        P_{n+1}=\min_{t\leq P_{n}!+1}(\text{Primo}(t) \land t>P_n)
      \end{equation}
      Questa restituisce il minimo numero primo maggiore di \(P_n\), ovvero il successivo numero primo. Infatti si ha:
      \begin{equation}
        \forall i\in\cbrakets{1,\cdots,n}, \neg(P_i|(P_n!+1))
      \end{equation}
      poiché\footnote{Si noti che \(P_n!+1=(1\cdot 2 \cdot \cdots \cdot P_i \cdot \cdots \cdot P_n)+1\), dunque \(\forall i\in\cbrakets{1,\cdots,n}: P_i|P_n!\), e dunque si ha che \((P_n!+1)\mod (P_i)=1\).}   \((P_n!+1)\mod (P_i)=1\), ovvero nessun primo \(\leq P_n\) divide \(P_n!+1\). Dunque qualsiasi fattore primo di \(P_n!+1\) sarà maggiore di \(P_n\), e in particolare\footnote{Questo limite superiore che imponiamo è molto più grande di quanto necessario, è dimostrabile che \(P_{n+1}\leq2P_n\). Si è scelto questo limite superiore essendo più facile da mostrare.}   si avrà \(P_n!+1\geq P_{n+1}\). % chktex 40

      Definiamo dunque:
      \begin{equation}
        h(y,z) := \min_{t\leq y}(\text{Primo}(t) \land t>y)
      \end{equation}
      questa è ricorsiva primitiva per il teorema di minimalizzazione limitata.

      Definiamo infine:
      \begin{equation}
        k(x):= h(x, x!+1)
      \end{equation}
      Questa è ricorsiva primitiva per composizione di funzioni ricorsive primitive. Possiamo infine definire per ricorsione primitiva \(P_n\) come:
      \begin{equation}
         \begin{cases}
            P_0=0\\
            P_{n+1}=k(P_n)
          \end{cases}
      \end{equation}
    \end{itemize}
  \end{enumerate}
\end{generalbox}

\newpage

\section{Funzioni parzialmente calcolabili}

Abbiamo analizzato le \textbf{classi PRC} nelle loro proprietà poiché abbiamo dimostrato che la classe delle \textbf{funzioni calcolabili è PRC}. \\
Dal capitolo~\ref{chapter 1} però, sappiamo che le \textbf{funzioni calcolabili} non sono le uniche di cui si può valutare la calcolabilità, infatti esse sono un sotto insieme delle \textbf{funzioni parzialmente calcolabili}.\\
Queste però sono escluse dalle caratterizzazioni delle \textbf{classi PRC} poiché definite solo per funzioni totali.\\
Per ovviare a questa lacuna possiamo fornirci di un ulteriore operatore, quello di \textbf{minimalizzazione illimitata} che, come vedremo, ci permetterà di caratterizzare a pieno la classe delle \textbf{funzioni parzialmente calcolabili}.

\subsection{Minimalizzazione illimitata}

\dfn{Minimalizzazione illimitata}{\label{Minimalizzazione illimitata Def}
  Sia \(p\) un predicato \((n+1)\text{-ario}\), allora chiamiamo \textbf{minimalizzazione illimitata} la funzione \((n+1)\text{-aria}\) \(\min_t\) definita come:
  \begin{equation}
      \min_t\; p(x_1,\cdots,x_n,t)= \begin{cases}
      \min\cbrakets{t\in\mbN | p(x_1,\cdots,x_n,t)} , & \exists t \in \mbN: p(x_1,\cdots,x_n,t)\\
      \uparrow, &\text{altrimenti}
    \end{cases}
  \end{equation}
}

Per far si che quest'operatore sia utile nella caratterizzazione delle \textbf{funzioni parzialmente calcolabili} è necessario dimostrare che esso stesso è \textbf{parzialmente calcolabile}.

\begin{halfframedbox}{red!75!black}{Calcolabilità della Minimalizzazione illimitata}\label{Calcolabilità Minimalizzazione illimitata}
  \textbf{Enunciato.} Se p un predicato calcolabile. Allora \(m\) \textbf{minimalizzazione illimitata} è parzialmente calcolabile.

  \begin{center}
    \textcolor{red!75!black}{\hdashrule[0.5ex]{18cm}{1pt}{3mm 3pt 1mm 2pt}}
  \end{center}

  \textbf{Dimostrazione.} Non potendo utilizzare le proprietà delle \textbf{classi PRC} ricorriamo al seguente S-programa:
\begin{lstlisting}[language=ini, label={Minimalizzazione illimitata Dim}, caption={Minimalizzazione illimitata}]
[A] IF (*\(p(X_1,\cdots,X_n,Y)\)*) GOTO E
    Y <- Y + 1
    GOTO A
\end{lstlisting}

  Si avrà infatti che se \(p(x_1,\cdots,x_n,t)\) è vera per qualche \(t\), il programma terminerà restituendo il minimo \(t\)  trovato, altrimenti il programma non terminerà mai.
\end{halfframedbox}

\subsection{Caratterizzazione delle funzioni parzialmente calcolabili}

Grazie all'operatore appena definito, possiamo infine caratterizzare le \textbf{funzioni parzialmente calcolabili} grazie a un teorema che, come vedremo, è molto importante per la teoria della calcolabilità.

\begin{halfframedbox}{red!75!black}{Caratterizzazione delle funzioni parzialmente calcolabili}\label{Caratterizzazione delle funzioni parzialmente calcolabili}
  \textbf{Enunciato.} Una funzione è parzialmente calcolabile \textit{se e solo se} si può ottenere dalle funzioni iniziali mediante un numero finito di \textbf{composizioni}, \textbf{ricorsioni primitive} (applicata solo a eventuali funzioni totali ottenute) e \textbf{minimalizzazioni illimitate}.
\end{halfframedbox}

Questo teorema, che non dimostreremo, insieme alla tesi di \textbf{Church-Turing}\footnote{Enunciata nella definizione~\ref{Tesi di Church-Turing Def}}, che ci dice che ogni \textbf{funzione parzialmente calcolabile} è esprimibile mediante un \textbf{S-programma}, esprime un risultato teorico molto importante, ovvero che \textbf{ogni procedura algoritmica} può essere \textbf{ridotto alle funzioni iniziali} mediante questi \textbf{tre operatori}. \\
Per analizzare meglio questi programmi però, e farli valutare da delle funzioni, sarà necessario che anch'essi possano essere associati univocamente con dei numeri\footnote{Esercizi Consigliati:~\ref{Funzione conta primi},~\ref{Funzione floor di sqrt2 x} e~\ref{Massimalizzazione limitata Eq}}.